\hypertarget{cluster-computing-with-pyspark-and-dask}{%
\chapter{Cluster Computing with PySpark and
Dask}\label{cluster-computing-with-pyspark-and-dask}}

\hypertarget{pyspark-the-jvm-bridge}{%
\subsection{78.1 PySpark: The JVM Bridge}\label{pyspark-the-jvm-bridge}}

PySpark is a Python wrapper for Apache Spark (written in Scala/JVM). *
\textbf{The Architecture}: Python code uses the \textbf{Py4J} bridge to
communicate with the Spark JVM. * \textbf{RDDs and DataFrames}: These
are distributed data structures that partitioned across the cluster. *
\textbf{Godhood Tip}: Avoid UDFs (User Defined Functions) in PySpark if
possible, as they require moving data between the JVM and Python
process, which is a massive performance bottleneck. Use Spark SQL
expressions instead.

\hypertarget{dask-native-python-parallelism}{%
\subsection{78.2 Dask: Native Python
Parallelism}\label{dask-native-python-parallelism}}

Unlike Spark, Dask is written entirely in Python. * \textbf{Task
Graphs}: Dask creates a DAG (Directed Acyclic Graph) of operations. *
\textbf{Schedulers}: Dask can run on a single machine (using
threads/processes) or on a distributed cluster of thousands of nodes.

\begin{center}\rule{0.5\linewidth}{0.5pt}\end{center}
